\chapter{The M/G/1 Queue and the Pollaczek--Khinchine Formula: the Price of a Miss}\label{sec:L3}

In \cref{sec:L2} the service times were exponential, so the state of the queue
was summarised by a single number, the number of customers. Prefill work is not
exponential. It is a general distribution that mixes short appends (hits) with
long recomputations (misses). The goal of this chapter is to compute the mean
waiting time of a FIFO queue with a general service distribution (the
Pollaczek--Khinchine formula), and then to differentiate it to prove the central
proposition of the paper, the price of a miss. The key message is a single one.
\emph{The mean wait of a FIFO queue depends not only on the mean of the work but
also on its second moment.}

\section{The M/G/1 queue and the residual service time}

The M/G/1 queue (\cref{L1:def:queues}(c)) has Poisson arrivals of rate $\lambda>0$,
i.i.d.\ work $S_k$ with an arbitrary law $G$ (``General''), independent of the
arrivals, with generic $S$ and $\E[S]<\infty$, and one FIFO server that serves each
customer to completion without interruption. The load is $\rho=\lambda\E[S]$
(\cref{def:load}). As in \cref{L1:def:measures}, $W_q$ is the wait from arrival to the
start of service, the sojourn time is $W_q+S$, and $L_q$ and $L$ are the time averages
of the number of customers in the waiting room and in the whole system.

Unlike the exponential distribution, a general service distribution is not
memoryless (\cref{prop:memoryless}). How much longer a customer who has already
been served for 3 seconds will take depends on the information ``3 seconds''.
The wait of an arriving customer therefore requires the quantity ``how much
service is left for the customer now in service''.

\begin{definition}[Residual service time]\label{L3:def:residual}
If the server is serving a customer at time $t$, let $R(t)$ be the remaining
service time of that customer; if the server is idle, set $R(t)=0$. $R(t)$ is
called the \term{residual service time}.
\end{definition}

The sample path of $R(t)$ is a sawtooth. At the moment service of customer $k$
begins, it jumps up to $S_k$; it then decreases with slope $-1$, and when it
reaches 0 the next service begins or the server becomes idle. Hence the
contribution of a single customer $k$ to $\int R(t)\dd t$ is the area $S_k^2/2$
of a triangle with base and height $S_k$.

\begin{lemma}[Time average of the residual service time \cite{kleinrock1975}]\label{lem:residual}
In an M/G/1 queue, suppose $\E[S^2]<\infty$ and that the number $D(T)$ of
customers who have completed service by time $T$ satisfies $D(T)/T\to\lambda$
(a.s.) (stability; this holds if $\rho<1$). Then almost surely
\[
  \lim_{T\to\infty}\frac1T\int_0^T R(t)\dd t=\frac{\lambda\E[S^2]}{2}
  =\rho\cdot\frac{\E[S^2]}{2\E[S]}.
\]
In particular, averaged only over the time the server is busy, the mean residual
service time is $\E[S^2]/(2\E[S])$.
\end{lemma}

\begin{proof}
Let $A(T)$ be the number of arrivals in $[0,T]$. The triangles of the customers
who completed service in $[0,T]$ lie entirely in $[0,T]$, and every customer
whose service has started arrived before $T$, so
\[
  \sum_{k=1}^{D(T)}\frac{S_k^2}{2}\ \le\ \int_0^T R(t)\dd t\ \le\ \sum_{k=1}^{A(T)}\frac{S_k^2}{2}.
\]
(Under FIFO the service order equals the arrival order, so the first $D(T)$
customers are exactly those who have completed service.) By the strong law of
large numbers $\frac1n\sum_{k\le n}S_k^2\to\E[S^2]$; for a Poisson process
$A(T)/T\to\lambda$, and by assumption $D(T)/T\to\lambda$. Hence
\[
  \frac1T\sum_{k=1}^{D(T)}\frac{S_k^2}{2}
  =\frac{D(T)}{T}\cdot\frac{1}{D(T)}\sum_{k=1}^{D(T)}\frac{S_k^2}{2}
  \longrightarrow\lambda\cdot\frac{\E[S^2]}{2},
\]
and the upper sum has the same limit. The squeeze gives the first equality. The
second equality is $\lambda=\rho/\E[S]$. The fraction of time the server is busy
is $\rho$ (\cref{cor:busy}) and $R=0$ during idle time, so the average over busy
time is $\lambda\E[S^2]/(2\rho)=\E[S^2]/(2\E[S])$.
\end{proof}

\begin{remark}[Inspection paradox]\label{L3:rem:inspection}
$\E[S^2]/(2\E[S])=\frac{\E[S]}{2}(1+\mathrm{CV}^2)$, where
$\mathrm{CV}^2=\Var[S]/\E[S]^2$ is the \term{squared coefficient of variation}.
The intuition ``a customer in service has, on average, about half of its service
left'' fails because, when we look at the server at a random time, the
probability of catching a long service is proportional to its length. Indeed,
the service observed at a busy instant has the \term{length-biased
distribution} with density $x\,G(\mathrm{d}x)/\E[S]$, whose mean is
$\E[S^2]/\E[S]$; the observation instant is uniform within that service, so the
mean remaining time is half of it. This is the \term{inspection paradox}. For the
exponential distribution $\mathrm{CV}^2=1$, so the mean residual equals $\E[S]$
itself, in agreement with memorylessness.
\end{remark}

\section{The Pollaczek--Khinchine formula}

\begin{theorem}[Pollaczek--Khinchine mean formula \cite{pollaczek1930,khinchine1932}]\label{thm:pk}
In an M/G/1 FIFO queue let $\rho<1$ and $\E[S^2]<\infty$. Then in steady state
\begin{equation}\label{L3:eq:pk}
  \E[W_q]=\frac{\lambda\E[S^2]}{2(1-\rho)},\qquad
  L=\rho+\frac{\lambda^2\E[S^2]}{2(1-\rho)}.
\end{equation}
\end{theorem}

\paragraph{Steady-state assumption.} The proof takes the following facts as
given. If $\rho<1$, the queue is stable, the waiting time $W_q^{(k)}$ of the
$k$-th customer converges in distribution, and the limit law has a finite mean
(when $\E[S^2]<\infty$). Moreover the customer averages and the time averages
exist almost surely. These are standard results obtained by analysing the
Lindley recursion $W^{(k+1)}_q=(W^{(k)}_q+S_k-\tau_{k+1})^+$ ($\tau$ the
interarrival times); see \cite[Ch.~3, 10]{asmussen2003} and
\cite[Ch.~5]{kleinrock1975}. The proof below accepts this existence and computes
the \emph{value} of the limit.

\begin{proof}
Consider a customer who has just arrived in steady state (the ``tagged
customer''). Under FIFO this customer waits for (1) the residual service time of
the customer in service and (2) the service times of the $N_q$ customers already
in the waiting room:
\[
  W_q=R+\sum_{j=1}^{N_q}S_{(j)},
\]
where $R$ is the residual service time at the arrival instant and $S_{(j)}$ is the
service time of the $j$-th waiting customer.

\emph{Mean of (1).} Poisson arrivals see time averages (PASTA, \cref{thm:pasta}).
So the law of $R$ seen at the arrival instant is the time-stationary law of
$R(t)$, and by \cref{lem:residual} $\E[R]=\lambda\E[S^2]/2$.

\emph{Mean of (2).} The waiting customers have not started service yet. Under
FIFO the state of the queue up to the arrival of the tagged customer is
determined only by the arrival times and the service times of the customers who
have already started service. Hence the service times $S_{(1)},S_{(2)},\dots$ of
the customers not yet started can be regarded as i.i.d.\ $G$ random variables
independent of $N_q$. Therefore
$\E\bigl[\sum_{j\le N_q}S_{(j)}\bigr]=\sum_{n}\Prob(N_q=n)\,n\E[S]=\E[N_q]\E[S]$.
By PASTA again, $\E[N_q]$ equals the time average $L_q$ of the number in the
waiting room.

\emph{Combining.} Applying Little's law (\cref{thm:little}) to the waiting room
alone as the ``system'' gives $L_q=\lambda\E[W_q]$. Therefore
\[
  \E[W_q]=\frac{\lambda\E[S^2]}{2}+\lambda\E[W_q]\,\E[S]
  =\frac{\lambda\E[S^2]}{2}+\rho\,\E[W_q].
\]
Since $\E[W_q]<\infty$, subtracting $\rho\E[W_q]$ from both sides and dividing by
$1-\rho>0$ gives the first formula. The second formula is Little's law applied to
the whole system: $L=\lambda(\E[W_q]+\E[S])=\rho+\lambda\E[W_q]$.
\end{proof}

\begin{corollary}[Variance increases the wait \cite{kleinrock1975}]\label{cor:variance}
Under the assumptions of \cref{thm:pk},
\begin{equation}\label{L3:eq:pk-cv}
  \E[W_q]=\frac{\rho}{1-\rho}\cdot\frac{1+\mathrm{CV}^2}{2}\cdot\E[S].
\end{equation}
Hence, of two work laws with the same $\lambda$ and the same mean $\E[S]$ (so the
same load $\rho<1$), the one with the larger variance waits strictly longer. In
particular:
\begin{enumerate}[label=(\roman*)]
\item M/M/1 ($\mathrm{CV}^2=1$): $\E[W_q]=\rho\E[S]/(1-\rho)$, in agreement with
  \cref{prop:mm1}.
\item M/D/1 (deterministic service, $\mathrm{CV}^2=0$):
  $\E[W_q]=\rho\E[S]/(2(1-\rho))$, exactly half of M/M/1 at the same load.
\end{enumerate}
\end{corollary}

\begin{proof}See \cref{L3:ex:variance}.\end{proof}

\begin{remark}
Contrast this with the PS queue of \cref{sec:L2} (\cref{def:ps}). The mean number
in an M/G/1-PS queue is $\rho/(1-\rho)$, insensitive to the work law
(\cref{thm:insens}), while the mean number in FIFO, \cref{L3:eq:pk}, is sensitive
to the law through $\E[S^2]$. A single long job blocks every customer behind it
under FIFO, but not under PS, where all customers share the server.
\end{remark}

\section{Mixtures of hits and misses}

The KV cache policy decides whether each turn is a hit or a miss, and hence sets
the law of the prefill work.

Let $S^{\mathrm{hit}}$ and $S^{\mathrm{miss}}$ be square-integrable nonnegative
random variables with $S^{\mathrm{miss}}\ge S^{\mathrm{hit}}$ (a.s.), and let $U$ be
uniform on $(0,1)$ and independent of them. At \term{hit rate} $h\in[0,1]$ the work is
the \term{hit/miss mixture}
\begin{equation}\label{L3:eq:mixture}
  S_h=\1\{U\le h\}\,S^{\mathrm{hit}}+\1\{U>h\}\,S^{\mathrm{miss}},
\end{equation}
which has the law of $S^{\mathrm{hit}}$ with probability $h$ and that of
$S^{\mathrm{miss}}$ with probability $1-h$.

\begin{corollary}[A higher hit rate lowers the PK wait]\label{cor:mixture}
Let $m_1(h)=\E[S_h]$, $m_2(h)=\E[S_h^2]$ and $\rho(h)=\lambda m_1(h)$.
\begin{enumerate}[label=(\roman*)]
\item $m_1$, $m_2$ and $\rho$ are affine and nonincreasing in $h$.
\item If $\rho(h_0)<1$, then $\rho(h)<1$ for all $h\in[h_0,1]$, and
  $W(h)=\lambda m_2(h)/(2(1-\rho(h)))$ is nonincreasing on $[h_0,1]$.
\end{enumerate}
\end{corollary}

\begin{proof}See \cref{L3:ex:mixture}.\end{proof}

\section{The price of a miss}

\Cref{cor:mixture} gives only a direction. A scheduler needs a \emph{magnitude}.
If we drop the KV state of a particular program $i$, so that its next turn is a
miss, by how much does the delay of the whole prefill stage grow?

Model the prefill stage as an M/G/1 FIFO queue with arrival rate $\lambda$, work
$S$ ($m_1=\E[S]$, $m_2=\E[S^2]<\infty$) and load $\rho=\lambda m_1<1$. Let
$W=\lambda m_2/(2(1-\rho))$ be the PK wait and $L_P$ the mean number in the
prefill stage (waiting and in service). By Little's law
\begin{equation}\label{L3:eq:LP}
  L_P=\lambda(W+m_1)=\rho+\frac{\lambda^2m_2}{2(1-\rho)}.
\end{equation}
For programs (classes) $i=1,\dots,k$, a \term{policy change} turns a fraction
$q_i\ge0$ of arrivals ($\sum_iq_i\le1$) from hits with work
$S_i^{\mathrm{hit}}\ge0$ into misses with work
$S_i^{\mathrm{miss}}=S_i^{\mathrm{hit}}+\Delta S_i$, $\Delta S_i\ge0$. The moments of
the work become
\[
  m_1'=m_1+\delta_1,\quad \delta_1=\sum_iq_i\Delta S_i,\qquad
  m_2'=m_2+\delta_2,\quad \delta_2=\sum_iq_i\bigl((S_i^{\mathrm{miss}})^2-(S_i^{\mathrm{hit}})^2\bigr),
\]
and the new load is $\rho'=\rho+\lambda\delta_1$. We write $L_P'$ for the new mean
number and $\Delta L_P=L_P'-L_P$ for the change.

\begin{theorem}[The price of a miss]\label{thm:price}
In this setting, suppose $\rho'<1$ and define
\begin{equation}\label{L3:eq:price}
  \Phi_i=\Delta S_i
  +\frac{\lambda\bigl((S_i^{\mathrm{miss}})^2-(S_i^{\mathrm{hit}})^2\bigr)}{2(1-\rho)}
  +\frac{\lambda W\,\Delta S_i}{1-\rho}.
\end{equation}
Then
\[
  \lambda\sum_iq_i\Phi_i\ \le\ \Delta L_P\ \le\ \frac{1-\rho}{1-\rho'}\,\lambda\sum_iq_i\Phi_i .
\]
\end{theorem}

\begin{proof}
Write $u=1-\rho>0$. Then $1-\rho'=u-\lambda\delta_1>0$, and since $\delta_1\ge0$,
$0<u-\lambda\delta_1\le u$. Applying the PK formula (\cref{thm:pk}) to the new work,
\[
  \Delta L_P=\lambda\delta_1+\frac{\lambda^2}{2}\Bigl(\frac{m_2+\delta_2}{u-\lambda\delta_1}-\frac{m_2}{u}\Bigr)
  =\lambda\delta_1+\frac{\lambda^2}{2}\cdot\frac{u(m_2+\delta_2)-(u-\lambda\delta_1)m_2}{u(u-\lambda\delta_1)}
  =\lambda\delta_1+\frac{Y}{u(u-\lambda\delta_1)},
\]
where $Y=\frac{\lambda^2}{2}\bigl(u\delta_2+\lambda\delta_1m_2\bigr)\ge0$. On the other hand,
taking the $q_i$-weighted sum of \cref{L3:eq:price} and using $W=\lambda m_2/(2u)$,
\[
  \lambda\sum_iq_i\Phi_i=\lambda\delta_1+\frac{\lambda^2\delta_2}{2u}+\frac{\lambda^2W\delta_1}{u}
  =\lambda\delta_1+\frac{\lambda^2}{2u^2}\bigl(u\delta_2+\lambda\delta_1m_2\bigr)
  =\lambda\delta_1+\frac{Y}{u^2}.
\]
\emph{Lower bound.} Since $u-\lambda\delta_1\le u$, $Y/(u(u-\lambda\delta_1))\ge Y/u^2$, and
hence $\Delta L_P\ge\lambda\sum_iq_i\Phi_i$.

\emph{Upper bound.} A direct computation gives
\[
  \frac{u}{u-\lambda\delta_1}\Bigl(\lambda\delta_1+\frac{Y}{u^2}\Bigr)-\Delta L_P
  =\lambda\delta_1\Bigl(\frac{u}{u-\lambda\delta_1}-1\Bigr)+\frac{Y}{u(u-\lambda\delta_1)}-\frac{Y}{u(u-\lambda\delta_1)}
  =\frac{(\lambda\delta_1)^2}{u-\lambda\delta_1}\ \ge0 .
\]
Since $u/(u-\lambda\delta_1)=(1-\rho)/(1-\rho')$, the upper bound holds.
\end{proof}

\paragraph{Reading the three terms.} $\Phi_i$ is measured in seconds. Misses of
program $i$ occur at rate $\lambda q_i$, so by Little's law $\lambda q_i\Phi_i$ is
an increase in the mean number, and $\Phi_i$ is ``the total waiting time, summed
over all turns, that one miss adds''.
\begin{enumerate}[label=(\arabic*)]
\item $\Delta S_i$: the miss itself takes longer to serve (its own cost).
\item $\lambda((S^{\mathrm{miss}}_i)^2-(S^{\mathrm{hit}}_i)^2)/(2(1-\rho))$: turns
  that arrive while the miss is being served wait for all of it. This comes from
  the second moment in the PK formula, and its size is the difference of the
  \emph{squares} of the work. It is the cost imposed on other customers, the
  \term{externality} of the economics of queues; in these notes we call it the
  \term{head-of-line term}.
\item $\lambda W\Delta S_i/(1-\rho)$: the extra work raises the load, which
  stretches every wait by the factor $1/(1-\rho)$ (the load term).
\end{enumerate}
Since $(S^{\mathrm{miss}})^2-(S^{\mathrm{hit}})^2=(S^{\mathrm{miss}}+S^{\mathrm{hit}})\Delta S$,
\[
  \frac{\text{(2)}}{\text{(1)}}=\frac{\lambda(S_i^{\mathrm{miss}}+S_i^{\mathrm{hit}})}{2(1-\rho)},\qquad
  \frac{\text{(2)}}{\text{(3)}}=\frac{S_i^{\mathrm{miss}}+S_i^{\mathrm{hit}}}{2W}.
\]
The first ratio is of the order of ``the number of turns that arrive while one
miss is served''. When it is small, the price is close to $\Delta S_i$; when it is
large, the head-of-line term dominates. Once the mean wait $W$ exceeds half the
sum of the miss and hit work, the load term exceeds the head-of-line term.

\begin{example}[At moderate load the head-of-line term dominates]\label{L3:ex:price}
In seconds, let $\lambda=1.25$, $m_1=0.4$, $m_2=0.5$, so $\rho=0.5$ and $W=0.625$.
Let a program have hit work $S^{\mathrm{hit}}=0.2$ and miss work
$S^{\mathrm{miss}}=2.0$ ($\Delta S=1.8$). The three terms are
\[
  \Delta S=1.8,\qquad \frac{1.25\,(4-0.04)}{2\cdot0.5}=4.95,\qquad
  \frac{1.25\cdot0.625\cdot1.8}{0.5}=2.8125,
\]
so $\Phi=9.5625$ seconds. One miss creates more than five times its own extra
service in total waiting, and most of it is paid by the turns queued behind it.
The two ratios are $2.75$ and $1.76$. If a fraction $q=0.01$ of this program's
turns turn into misses, then $\rho'=0.5225$, and the exact change
$\Delta L_P\approx0.1241$ lies in the interval $[0.1195,\ 0.1252]$ of the theorem.
\end{example}

\begin{corollary}[Divergence near saturation]\label{cor:price-diverge}
Fix the work law ($m_1>0$, $m_2>0$) and $\Delta S_i>0$, and let
$\lambda\uparrow1/m_1$ (that is, $\rho\uparrow1$). Then $\Phi_i\to\infty$.
\end{corollary}

\begin{proof}See \cref{L3:ex:diverge}.\end{proof}

\begin{corollary}[Monotonicity in the context length]\label{cor:price-mono}
Let the prefill cost be $P(n,K)=an+bn(K+n/2)$ ($a,b\ge0$; the cost of computing
$n$ tokens on top of a prefix of length $K$). A hit costs
$S^{\mathrm{hit}}=P(n,K)$ and a miss, which recomputes the whole context, costs
$S^{\mathrm{miss}}=P(K+n,0)$. Then $\Delta S=aK+bK^2/2$, and with $\lambda,\rho,W$
and the number $n$ of appended tokens fixed, $\Phi$ is nondecreasing in $K\ge0$.
\end{corollary}

\begin{proof}See \cref{L3:ex:mono}.\end{proof}

\serving{Chunked prefill interleaves chunks of a long prefill with decode steps, but
pending prefills are still served in arrival order: a prefill behind a long miss waits
for every chunk of it, as long as the oldest prefill may take all the budget the
decode batch leaves. Then the head-of-line term (2) remains, whatever the chunk size; a
per-prefill cap lets later prefills bypass the miss (\cref{L5:ex:colocated}).
\Cref{cor:price-mono} says a longer context has a larger price, but the scheduler
compares $p_i\Phi_i/c_i$, and that order can flip with the load (\cref{L5:ex:flip}).}

\begin{remark}[Limits of the M/G/1 assumption]\label{L3:rem:finite}
A pool holds a few dozen live sessions, and a session waiting for prefill cannot create
its next turn, so arrivals slow down when the queue is long. For exponential work the
finite-population model of \cref{sec:L4} shows that the open wait is an upper bound
(\cref{prop:finite}); for general work the comparison is a question for measurement.
\end{remark}

\section{Summary}
\begin{itemize}
\item The FIFO wait depends on the second moment of the work, $\E[W_q]=\lambda\E[S^2]/(2(1-\rho))$.
\item The price of a miss has an own-work, a head-of-line and a load term, and brackets the change in the mean number at prefill.
\item It diverges near saturation and grows with the context; chunked prefill does not remove the head-of-line term.
\end{itemize}

\section{Exercises}

\begin{exercise}\label{L3:ex:variance}
Derive \eqref{L3:eq:pk-cv} from \cref{thm:pk} and check cases (i) and (ii) of
\cref{cor:variance}.
\emph{Hint:} $\E[S^2]=\E[S]^2(1+\mathrm{CV}^2)$ and $\lambda\E[S]=\rho$.
\end{exercise}

\begin{exercise}\label{L3:ex:mixture}
Prove \cref{cor:mixture}.
\emph{Hint:} $m_k(h)$ is affine in $h$ with slope
$\E[(S^{\mathrm{hit}})^k]-\E[(S^{\mathrm{miss}})^k]\le0$; in (ii) the numerator of $W$
falls and its denominator rises.
\end{exercise}

\begin{exercise}\label{L3:ex:diverge}
Prove \cref{cor:price-diverge}.
\emph{Hint:} with $\rho=1-\varepsilon$ and $\varepsilon\le1/2$, the load term of
\eqref{L3:eq:price} is $\lambda^2m_2\Delta S_i/(2\varepsilon^2)\ge
m_2\Delta S_i/(8m_1^2\varepsilon^2)$.
\end{exercise}

\begin{exercise}\label{L3:ex:mono}
Prove \cref{cor:price-mono}.
\emph{Hint:} expand $P(K+n,0)-P(n,K)$; each term of \eqref{L3:eq:price} is a
nonnegative constant times $\Delta S$, $(S^{\mathrm{miss}}+S^{\mathrm{hit}})\Delta S$ or
$\Delta S$, and $S^{\mathrm{miss}}+S^{\mathrm{hit}}=2S^{\mathrm{hit}}+\Delta S$.
\end{exercise}

\begin{exercise}\label{L3:ex:1}
Show that the mean residual service time over busy time, $\E[S^2]/(2\E[S])$, is
always at least $\E[S]/2$, with equality only when $S$ is deterministic. Also
check that it equals $\E[S]$ when $S$ is exponential.
\emph{Hint:} $\E[S^2]=\E[S]^2+\Var[S]$.
\end{exercise}

\begin{exercise}\label{L3:ex:2}
With the numbers of \cref{L3:ex:price} and $q=0.01$, compute $\Delta L_P$ exactly
and check $0.1195\le\Delta L_P\le0.1252$. How does the ratio of the upper to the
lower bound change as $q$ grows?
\emph{Hint:} substitute $m_1+\delta_1$ and $m_2+\delta_2$ into
$L_P=\lambda m_1+\lambda^2m_2/(2(1-\lambda m_1))$.
\end{exercise}
